\chapter*{Abstract}
\addcontentsline{toc}{chapter}{Abstract}

A programming tutor needs evidence to judge whether someone can solve the next problem. Asking for a short piece of code
and running it against tests offers one way to gather that evidence. Yet a test can miss a requirement, and its result can
be scored or interpreted incorrectly. This dissertation asks when an additional coding check helps predict success on a
separate task, and how unreliable checks affect that prediction.

The project builds an evaluation pipeline that records predictions before collecting the task outcomes used to assess them.
The first study uses answers and code from the language model \texttt{gpt-oss-120b} on 24 authored Python cases, with 23
available for the final comparison. A fault in the project's test software marked matching outputs as incorrect, affecting
both the evidence and the outcomes used to judge predictions. Repairing it reduced the apparent benefit of extra checks
from about 30 to 4 percentage points. Correct predictions increased from 18 to 19 out of 23 after correction.
The 95\% interval for the gain was 0 to 13 percentage points, leaving the size of any useful benefit uncertain.

Two computer simulations then examine how to handle and select checks when their reliability varies. Limiting how far one
check could change a probability estimate improved the accuracy of those estimates on average across the specified adverse conditions, although
some conditions became worse. A second study allowed checks for half of 40 cases in each simulated round, with 2,000 rounds
in each of seven settings. The proposed selection rule accounted for uncertainty about check reliability. It beat random
selection but lost to simpler alternatives, including choosing the cases with the most uncertain predictions. Mathematical
inspection found that its score could assign value to a check even when either result would leave the prediction unchanged.

To examine uncertainty using student data, an exploratory study analysed Python activity records from the 2019 Computer
Science Educational Data Mining data challenge. Among 729 predictions for 86 learners, selecting 361 uncertain predictions
captured 72.3\% of prediction errors, compared with 49.5\% expected from random selection. Whether reviewing those errors
would improve learning was not tested.

The contribution is a reproducible evaluation tool and an investigation of how faulty checks and decision rules can distort
performance predictions. The findings show why evidence quality needs its own evaluation before it informs tutoring decisions.
Improved student learning remains a question for a subsequent teaching study.
